\section{References}
\begingroup\small
\begin{itemize}[leftmargin=1.25em,itemsep=0.2em]
\item Sacred Scripture, \emph{Douay--Rheims Bible}, Challoner revision, American edition (1899), retained eBible text: Genesis 2:17; 3:19; Matthew 6:19--21; 10:26--33; Mark 14:32--42; Luke 12:15--21, 35--40; 13:1--5; 22:54--62; 23:32--49; John 6:55; 10:17--18; 11:17--44; 19:25--30; 20:19--29; Romans 5:12--21; 6:3--13, 23; 8:31--39; 14:7--9; 1 Corinthians 15:12--49; Philippians 1:20--26; 1 Thessalonians 4:12--17; Hebrews 2:1--18; 4:14--16; 5:7--9; James 5:14--15. John 6:55 corresponds to John 6:54 in CCC 1524; the Douay--Rheims 1 Thessalonians 4:12--17 corresponds to 4:13--18 in many modern editions. The brief Romans 14:8 quotation preserves the identified historical wording.
\item \emph{Catechism of the Catholic Church}, official English text, Holy See: 624--630 (Christ's real death and continuing personal union with body and soul); 988--1014 (resurrection, dying in Christ, and preparation); 1021--1032 (judgment, heaven, purification); 1514--1525 (anointing and Viaticum); 1684--1690 (funerals); 2276--2279 (care of the dying). The citation of Benedict XII's \emph{Benedictus Deus} follows the source reproduced at CCC 1023. Official section pages consulted 13 October 2026; current pages checked for anointing and care.
\item Second Vatican Council, \href{https://www.vatican.va/archive/hist_councils/ii_vatican_council/documents/vat-ii_const_19651207_gaudium-et-spes_en.html}{\emph{Gaudium et spes}} (7 December 1965), no.~18. Official English text: the human wound of death and Christ's victory.
\item Congregation for the Doctrine of the Faith, \href{https://www.vatican.va/roman_curia/congregations/cfaith/documents/rc_con_cfaith_doc_19790517_escatologia_en.html}{\emph{Letter on Certain Questions Concerning Eschatology}} (17 May 1979), nos.~2--5 and concluding doctrinal paragraphs. Official English text: the whole person, the subsisting human self, the Church's prayer, future resurrection, and the limits of imaginative descriptions.
\item St.~Augustine, \emph{The City of God}, XIII.3--6, especially 4--5, trans. Marcus Dods, vol.~I (Edinburgh: T.\ \& T.\ Clark, 1871), Project Gutenberg eBook 45304, retained transcription: why the baptized still die and how grace enables a good use of death.
\item St.~Augustine, \emph{Confessions}, IX.10.23--13.37, trans. J.~G. Pilkington, in \emph{Nicene and Post-Nicene Fathers}, first series, vol.~1, retained complete CCEL electronic delivery: contemplation at Ostia, Monica's death, Augustine's tears, and prayer at the altar. The translation and the later editorial notes are distinct; the essay draws on Augustine's text.
\item St.~Cyprian, \emph{On Mortality} (\emph{De mortalitate}), nos.~3--8, 16--18, 20--21, 26, in \emph{Ante-Nicene Fathers}, vol.~5, trans. Robert Ernest Wallis, retained complete CCEL delivery: shared bodily mortality, the struggle with vice, charity in an epidemic, martyrdom's dependence on God's gift, and desire for the heavenly homeland.
\item Church of Smyrna, \emph{The Martyrdom of Polycarp}, chapters 1, 4--5, 7--12, 14, 17--18, trans. Alexander Roberts and James Donaldson, in \emph{Ante-Nicene Fathers}, vol.~1, retained complete CCEL delivery: patient fidelity, hospitality and intercession, grateful confession, bodily resurrection, love of the martyrs, and annual commemoration as preparation. The complete constituent work, chapters 1--22 with apparatus, was read; editorial additions are distinguished from the ancient narrative.
\item St.~Ambrose, \emph{On the Death of His Brother Satyrus} (\emph{De excessu fratris Satyri}), II.1--3, 8; book~II is titled \emph{On the Belief in the Resurrection} in the English edition. Trans. H.~de Romestin, with E.~de Romestin and H.~T.~F. Duckworth, \emph{Nicene and Post-Nicene Fathers}, second series, vol.~10, retained complete CCEL delivery: grief, consolation, and grateful remembrance.
\item St.~Benedict, \emph{Rule}, chapter~4, Latin Wikisource witness, retained normalized text: the tools of good works, daily remembrance of death, works of mercy, reconciliation, and refusal of despair. The essay paraphrases; it does not offer a new translation for prayer.
\item St.~Francis de Sales, \emph{Introduction to the Devout Life: A New Translation} (London, Oxford, and Cambridge: Rivingtons, 1876), I.13, pp.~35--38. Translator unnamed in the inspected edition. Retained historical OCR checked against the printed page images: uncertain hour, universal farewell, holy friendship, and preparation of conscience.
\item St.~Alphonsus Liguori, \emph{Preparation for Death; or, Considerations on the Eternal Maxims}, trans. ``a Catholic clergyman'' (Dublin: James Duffy, 1843), IV.1 and V.1--2, especially pp.~60--61, 70--71, 74--75. The printed title date controls over the host's 1893 catalog label. Retained historical OCR and checked page images: certain death, uncertain hour, and the danger of postponing conversion. Quotations of other authors embedded in Alphonsus are not presented as independently checked witnesses.
\item St.~Thomas Aquinas, \emph{Summa theologiae} III, q.~50, aa.~2--3, trans. Fathers of the English Dominican Province, Project Gutenberg eBook 19950 (2006 electronic edition): the Word's continuing personal union with Christ's human soul and body in death.
\end{itemize}
The historical Scripture and translations retain their public-domain status in the United States; modern Holy See texts retain their own rights. Electronic transcriptions are not critical editions or facsimiles. No New Advent modernization or newly composed vocal prayer is reproduced.
\endgroup
